% !TeX root = ../main.tex
% Kompakte Beispielseite; verwendet die Pakete aus vorlage/praeambel.tex.
\section{LaTeX-Beispiele}
\label{sec:beispiele}
\begingroup
\small

\subsection{Text, Listen und Verweise}
Text kann \textbf{fett}, \textit{kursiv}, \underline{unterstrichen},
\texttt{in Schreibmaschinenschrift} oder \textcolor{Akzent}{farbig} sein.
Sonderzeichen: \verb|\% \& \_ \#| ergeben \%, \&, \_, \#.
Eine Fußnote\footnote{Hier steht eine ergänzende Erklärung.} und ein
\href{https://www.latex-project.org/}{Link zum LaTeX-Projekt} ergänzen den Text.

\noindent
\begin{minipage}[t]{0.48\linewidth}
  \begin{itemize}
    \setlength{\itemsep}{0pt}
    \setlength{\parskip}{0pt}
    \item Stichpunkt mit \emph{Betonung}
    \item Messwert: \qty{12.5}{\volt}
  \end{itemize}
\end{minipage}\hfill
\begin{minipage}[t]{0.48\linewidth}
  \begin{enumerate}
    \setlength{\itemsep}{0pt}
    \setlength{\parskip}{0pt}
    \item Versuchsaufbau prüfen
    \item Messung durchführen
  \end{enumerate}
\end{minipage}

\subsection{Formeln und Einheiten}
Im Fließtext gilt $U=R\cdot I$; Indizes und Potenzen schreibt man als
$U_{\mathrm{eff}}$ und $x^2$. Eine nummerierte Gleichung:
\begin{equation}
  R=\frac{U}{I},\qquad
  \overline{x}=\frac{1}{n}\sum_{i=1}^{n}x_i,\qquad
  \sigma=\sqrt{\frac{1}{n}\sum_{i=1}^{n}(x_i-\overline{x})^2}.
  \label{eq:auswertung}
\end{equation}
Einheiten und Unsicherheiten: $R=\qty{100.0 +- 0.5}{\ohm}$,
$I=\qty{25}{\milli\ampere}$ und $T=\qty{23}{\degreeCelsius}$.

\subsection{Tabelle und Abbildung}
% Ohne Gleitumgebungen bleiben beide Beispiele an dieser Stelle.
\noindent
\begin{minipage}[t]{0.48\linewidth}
  \centering
  \captionof{table}{Beispielmesswerte}
  \label{tab:messwerte}
  \begin{tabular}{@{}rSS@{}}
    \toprule
    Nr. & {$U$ in \unit{\volt}} & {$I$ in \unit{\milli\ampere}} \\
    \midrule
    1 & 1.0 & 10.0 \\
    2 & 2.0 & 20.1 \\
    3 & 3.0 & 29.9 \\
    \bottomrule
  \end{tabular}
\end{minipage}\hfill
\begin{minipage}[t]{0.48\linewidth}
  \centering
  \vspace{0pt}
  \fbox{\parbox[c][17mm][c]{0.85\linewidth}{\centering
    Platzhalter für ein Bild\\[2pt]
    \footnotesize\texttt{\string\includegraphics[width=\string\linewidth]}\\
    \texttt{\{dateiname\}}}}
  \captionof{figure}{Beispiel einer Bildunterschrift}
  \label{fig:beispiel}
\end{minipage}

Automatische Querverweise: \cref{eq:auswertung,tab:messwerte,fig:beispiel}.
Mit \verb|\label{...}| wird eine Marke gesetzt, mit \verb|\cref{...}| darauf verwiesen.

\subsection{Quelltext}
\begin{verbatim}
\textbf{Fetter Text}     % Kommentar im Quelltext
\qty{4.7}{\kilo\ohm}    % Zahlenwert mit Einheit
\end{verbatim}
\endgroup
